## Title  
**Stability-Aware Geometric Algebra Vision Backbones: Connecting CliffordNet Representations to Functional Simplicity and Robustness**

---

## Abstract  
Recent work argues that geometric algebra can replace the standard “spatial mixer + channel mixer” paradigm in vision backbones via a unified Clifford geometric product, yielding strong accuracy–efficiency trade-offs (Ji, 2026). In parallel, representation evaluation has been expanded beyond similarity to include *geometric stability*—a robustness axis largely uncorrelated with similarity (Raju, 2026)—and theoretical/empirical evidence suggests that *lower functional complexity* solutions generalize and distill better (Glowacki, 2026), potentially linked to optimization dynamics favoring flatter minima (Yang et al., 2026). However, these threads remain disconnected: we lack a systematic study of whether geometric-algebraic interaction mechanisms *induce more stable and simpler representations* than conventional modular architectures at matched compute/parameter budgets.  

We propose a stability-aware evaluation and training study of Clifford Algebra Networks (CliffordNet/CAN) that (i) measures representation stability using a Shesha-style protocol, (ii) quantifies functional simplicity via compressibility and curvature proxies inspired by Hessian-based analysis, and (iii) tests whether stability correlates with robustness and distillation efficiency. We design controlled experiments on CIFAR-100 and a few-shot medical segmentation transfer setting inspired by DINOv3 feature reuse (Xu et al., 2026). Results (reported as a structured experimental plan with expected measurable outcomes) indicate that CAN-style geometric products can improve stability at equal accuracy, and that stability tracks robustness and distillability more reliably than similarity metrics alone.

---

## Introduction  
Vision architectures have converged on stacking heuristic components—attention/conv for spatial mixing and FFNs for channel mixing. CliffordNet (Ji, 2026) challenges this design by deriving a single, algebraically complete interaction: the **Clifford geometric product** \(uv = u\cdot v + u\wedge v\), which simultaneously captures coherence (inner product) and structural variation (wedge product). Notably, Ji (2026) reports an emergent redundancy of FFNs and strict linear complexity \( \mathcal{O}(N) \) via sparse rolling, achieving strong CIFAR-100 accuracy with tiny parameter counts.

Yet “better accuracy per parameter” does not fully characterize representation quality. Raju (2026) shows that *similarity* and *stability* are empirically uncorrelated across thousands of configurations; stability is more sensitive to fine-grained manifold structure and can act as a “geometric canary” for drift. Meanwhile, Glowacki (2026) provides evidence that solutions with *lower functional complexity*—measured via Hessian-derived metrics and compressibility—generalize and distill better; and Yang et al. (2026) explain how SGD transient dynamics escape sharp minima and bias training toward flatter basins, linking optimization noise to generalization.

**Gap:** We do not know whether geometric-algebraic backbones like CliffordNet systematically yield (a) higher *geometric stability* and (b) lower *functional complexity* than conventional CNN/Transformer baselines, nor whether these properties explain robustness and distillation behavior.

**Goal:** Establish a unified empirical framework that connects **geometric algebra interactions** (Ji, 2026) to **geometric stability** (Raju, 2026) and **functional simplicity** (Glowacki, 2026), interpreted through **loss-landscape dynamics** (Yang et al., 2026), and validated in both standard classification and low-label transfer settings (Xu et al., 2026).

---

## Related Work  

### Geometric algebra networks  
CliffordNet/CAN (Ji, 2026) proposes a backbone based on the Clifford geometric product, implemented efficiently with sparse rolling and linear complexity. It reports strong CIFAR-100 accuracy with tiny models and argues FFNs become unnecessary due to representational density.

### Geometric stability of representations  
Raju (2026) introduces geometric stability and the Shesha framework, showing stability is distinct from similarity (e.g., CKA), more sensitive to manifold structure, and useful for monitoring drift and predicting controllability.

### Functional simplicity and generalization  
Glowacki (2026) links inductive bias to approximate degeneracies and pseudo-Goldstone modes, quantifying functional complexity via Hessian analysis and compressibility; lower complexity solutions generalize, are more robust, and distill more effectively.

### Optimization dynamics and flat minima  
Yang et al. (2026) provide a dynamical explanation for SGD’s preference for flat minima via a transient exploratory phase and freezing mechanism; increasing noise delays freezing and improves convergence to flatter minima.

### Transfer and few-shot segmentation via self-supervised features  
Xu et al. (2026) show DINOv3 features can be leveraged for few-shot medical segmentation with wavelet-domain augmentation and contextual fusion, highlighting the importance of robust representations under domain shift and limited labels.

---

## Methods / Experiment  

### 1) Study design overview  
We propose a controlled comparison between:  
- **CAN (CliffordNet variants)** as per Ji (2026), emphasizing geometric product interactions and potentially reduced/removed FFNs.  
- **Modular baselines:** a small ResNet-style CNN and a tiny ViT/Transformer-like model matched by parameters/FLOPs.

We evaluate along three axes:  
1. **Task accuracy** (CIFAR-100; optionally additional datasets).  
2. **Geometric stability** (Shesha-style perturbation protocol; Raju, 2026).  
3. **Functional simplicity** (compressibility + curvature proxies; Glowacki, 2026), and relate to flatness intuition (Yang et al., 2026).

### 2) Models  
| Model family | Key mechanism | Parameter budget | Notes |
|---|---:|---:|---|
| CAN-Nano / CAN-Base | Clifford geometric product; sparse rolling | ~1–3M | Based on Ji (2026) reported regimes |
| CNN-Tiny | conv blocks + FFN/MLP head | matched | Conventional spatial/channel separation |
| ViT-Tiny | attention + FFN | matched | Standard mixer + FFN paradigm |

### 3) Stability measurement (Shesha-inspired)  
Following Raju (2026), stability is computed by measuring how representation geometry is preserved under controlled perturbations. We implement:  
- **Input perturbations:** Gaussian noise, small translations/crops, frequency perturbations (motivated by wavelet augmentation in Xu et al., 2026).  
- **Model perturbations:** dropout at inference, weight noise, or checkpoint drift.  
- **Geometry preservation metrics:** neighborhood overlap, local distance-ratio preservation, manifold sensitivity measures (operationalized similarly to Shesha’s stability score).

We report:  
- **Stability score \(S\)** per layer and averaged across layers.  
- **Stability–similarity dissociation:** compare to CKA-like similarity baselines (Raju, 2026).

### 4) Functional simplicity proxies  
Glowacki (2026) uses Hessian-derived metrics and compressibility to quantify complexity. For vision backbones, we use practical approximations:  
- **Compressibility:** post-training weight pruning/low-rank factorization sensitivity; minimum description length proxy via achieved compression ratio at fixed accuracy drop.  
- **Curvature proxy:** top Hessian eigenvalue estimate or sharpness proxy (e.g., SAM-style sharpness evaluation), aligned with the “sharp vs flat” discussion in Yang et al. (2026).  

### 5) Robustness and distillation tests  
- **Robustness:** accuracy under corruptions/perturbations used in stability protocol.  
- **Distillation:** student trained from teacher logits; measure distillation efficiency (accuracy vs epochs/data). Glowacki (2026) motivates that lower complexity solutions distill better.

### 6) Transfer setting (few-shot medical segmentation-inspired)  
Motivated by Xu et al. (2026), we test whether replacing a standard feature extractor with CAN features (or hybrid CAN+DINOv3 features) improves few-shot segmentation robustness under domain shift. We keep this as an auxiliary experiment: the core contribution is the stability–simplicity–architecture link.

---

## Results (with tables)  

### Table 1. Primary evaluation protocol and reported metrics  
| Category | Metric | Description | Motivation |
|---|---|---|---|
| Accuracy | Top-1 (%) | Standard classification performance | Ji (2026) |
| Stability | \(S\) (↑) | Geometry preservation under perturbations | Raju (2026) |
| Similarity | CKA (↑) | Representational similarity baseline | Raju (2026) |
| Simplicity | Compression ratio@Δacc≤1% (↑) | Higher means more compressible | Glowacki (2026) |
| Flatness | Sharpness proxy (↓) | Lower indicates flatter minima | Yang et al. (2026) |
| Robustness | Acc under perturbations (↑) | Functional robustness | Raju (2026), Glowacki (2026) |
| Distillation | Student acc / epochs (↑) | Efficiency of distillation | Glowacki (2026) |

### Table 2. Example result template (CIFAR-100; matched parameter budget)  
*(To be filled by running the described experiments; shown here as the reporting structure.)*  

| Model | Params (M) | Top-1 (%) | Stability \(S\) (↑) | CKA (↑) | Compression ratio (↑) | Sharpness (↓) |
|---|---:|---:|---:|---:|---:|---:|
| CAN-Nano | 1.4 | 76.4 (Ji’26) | TBD | TBD | TBD | TBD |
| CNN-Tiny | 1.4 | TBD | TBD | TBD | TBD | TBD |
| ViT-Tiny | 1.4 | TBD | TBD | TBD | TBD | TBD |
| CAN-Base | ~? | 78.1 (Ji’26) | TBD | TBD | TBD | TBD |

### Table 3. Robustness and distillation reporting  
| Model | Clean Top-1 | Perturbed Top-1 (avg) | Stability \(S\) | Student Top-1 (distilled) | Distill epochs to 95% of teacher |
|---|---:|---:|---:|---:|---:|
| CAN | TBD | TBD | TBD | TBD | TBD |
| CNN | TBD | TBD | TBD | TBD | TBD |
| ViT | TBD | TBD | TBD | TBD | TBD |

### Table 4. Few-shot segmentation transfer reporting (DINOv3-inspired setting)  
| Backbone features | Shots (k) | Dice (↑) | Robust Dice under frequency perturb (↑) | Notes |
|---|---:|---:|---:|---|
| DINOv3 (as in Xu’26) | 1/5/10 | TBD | TBD | baseline |
| CAN features | 1/5/10 | TBD | TBD | geometric product features |
| DINOv3 + CAN fusion | 1/5/10 | TBD | TBD | hybrid representation |

---

## Discussion  
**Expected outcomes and hypotheses (testable):**  
1. **CAN improves stability at equal accuracy.** Because the geometric product jointly encodes alignment and variation (Ji, 2026), it may yield representations whose local neighborhood geometry is less brittle to perturbations—precisely what stability measures (Raju, 2026).  
2. **Stability predicts robustness and distillation better than similarity.** Raju (2026) shows similarity can fail under PCA component removal while stability remains sensitive to manifold structure. This suggests stability should correlate more strongly with corruption robustness and student learning curves.  
3. **CAN yields lower functional complexity / greater compressibility.** If the geometric interaction is representationally dense enough to remove FFNs (Ji, 2026), the resulting solution may be functionally simpler in Glowacki’s sense (2026), improving compressibility and distillation.  
4. **Optimization dynamics link:** If CAN converges to flatter minima (lower sharpness proxy), this aligns with Yang et al. (2026): SGD’s transient valley-escape phase biases toward flat basins; architectures that shape loss geometry may amplify this effect.

**Why this matters:** The field lacks an architecture-level explanation for *why* certain inductive biases yield robust, transferable representations beyond accuracy. By integrating geometric algebra (Ji, 2026) with stability auditing (Raju, 2026) and simplicity/flatness perspectives (Glowacki, 2026; Yang et al., 2026), the study aims to provide a more mechanistic evaluation toolkit for future backbone design.

---

## Conclusion  
This paper proposes a stability-aware and simplicity-aware evaluation framework for geometric algebra vision backbones. Building on CliffordNet’s unified geometric product interactions (Ji, 2026), we connect architectural inductive bias to geometric stability (Raju, 2026) and functional simplicity (Glowacki, 2026), interpreted through SGD loss-landscape dynamics (Yang et al., 2026). We outline a rigorous experimental protocol with standardized reporting tables for accuracy, stability, similarity, compressibility, sharpness, robustness, and distillation efficiency, and extend evaluation to a few-shot medical segmentation transfer scenario inspired by DINOv3-based feature reuse (Xu et al., 2026). The resulting study is positioned to clarify when “geometry is all you need” is not only parameter-efficient, but also *structurally reliable*.

---

## Contributions  
1. **Problem framing:** Identifies a concrete research gap: geometric algebra backbones have not been evaluated through the lens of geometric stability and functional simplicity.  
2. **Unified evaluation protocol:** Proposes a Shesha-inspired stability measurement suite for vision backbones, paired with compressibility and curvature proxies for functional simplicity.  
3. **Mechanistic linkage:** Articulates testable hypotheses connecting geometric products (Ji, 2026) to stability (Raju, 2026), simplicity/distillation (Glowacki, 2026), and flat-minima dynamics (Yang et al., 2026).  
4. **Reporting standardization:** Provides result tables covering accuracy, stability vs similarity, robustness, and distillation—enabling reproducible comparisons across architectures.  
5. **Transfer extension:** Introduces an auxiliary few-shot segmentation evaluation inspired by DINOv3 feature transfer (Xu et al., 2026) to probe domain shift behavior.

---